\section{The Level Predictability Score}\label{sec:lps}

The theory of \Cref{sec:theory} locates the \inm/\cnm{} crossover in terms of population quantities that are unobservable before training. This section operationalizes them as a diagnostic that can be computed from raw data and covariates alone, \emph{before} any forecasting model is fit.

\subsection{Definition}

Partition the training series into non-overlapping windows of length $w$ (stride $w$, so that cross-validation blocks are serially decoupled), and let $\bar{y}_i$ and $\bar{x}_i$ denote the target mean and the covariate mean of window $i$. The Level Predictability Score is the pooled out-of-sample $R^2$ of regressing window-level target means on window-level covariate means under expanding chronological cross-validation:
\begin{equation}\label{eq:lpsdef}
\lps \;=\; 1 \;-\; \frac{\sum_{k=1}^{K}\sum_{i\in\mathcal{T}_k}\bigl(\bar{y}_i - \hat{g}_k(\bar{x}_i)\bigr)^2}{\sum_{k=1}^{K}\sum_{i\in\mathcal{T}_k}\bigl(\bar{y}_i - \bar{y}^{\,\mathrm{tr}}_{k}\bigr)^2},
\end{equation}
where $\hat{g}_k$ is a first-stage regressor fit on all windows preceding test block $\mathcal{T}_k$, and the baseline $\bar{y}^{\,\mathrm{tr}}_{k}$ is the \emph{fold-train} mean of $\bar{y}$---never the out-of-sample mean, which would leak test-period level information into the denominator. We use $K=5$ folds with a minimum training fraction of $0.4$. The official instantiation fixes $w=96$ and takes $\hat{g}_k$ to be LightGBM \citep{ke2017lightgbm}, the same first-stage family used by CondNorm, so that $\lps$ measures exactly the information the conditional normalizer will have access to. For multivariate datasets, we average per-channel $\lps$ values, matching the channel-averaged test MSE of the experimental grid. The covariate set is identical to what CondNorm receives: lead-matched archived numerical weather forecasts and calendar harmonics for the exogenous group, calendar harmonics only for the standard benchmarks. A ridge first stage is reported as a sensitivity check but plays no role in the decision rule.

\subsection{Decision rule and threshold}

The diagnostic induces a one-line prescription, which we state in the \emph{negative} form that our evidence supports:
\begin{equation*}
\lps \geq \tau \;\Rightarrow\; \text{do not estimate the level endogenously};
\quad \text{otherwise instance normalization is the safe default},
\end{equation*}
with $\tau = 0.3$ fixed by pre-registration. The positive branch deliberately names no method. What \Cref{sec:empirical} establishes above the threshold is that \emph{several} covariate-conditioned level models---a first-stage regression used alone, a classical dynamic regression, and \cnm---all beat the best instance-normalized arm, and that they do not order among themselves in a stable way; a rule that named one of them would be claiming more than the data supports, and would be wrong on at least one dataset in our own panel. Below the threshold the prescription is directional in the other sense: every covariate-conditioned model we ran loses, so the endogenous default should be retained.

The value $\tau = 0.3$ was fixed in advance from the stylized model of \Cref{sec:theory}: under the parameter mapping of the synthetic study (\Cref{sec:synthetic}), the restoration-rule model M1 places its crossover \Cref{eq:lamstar} at $\lamstar \approx 0.27$--$0.28$ at $h=24$, and because M1 is an \emph{upper bound} on the crossover realized by trained predictors---whose empirical crossings in the synthetic grid occur at $0.01$--$0.03$ (\Cref{sec:synthetic})---adopting $\tau = 0.3$ is conservative toward instance normalization \citep{kim2022revin}. Two qualifications keep this from being read as more than it is. First, $\lps$ is not a consistent estimator of M1's $\lambda$: an out-of-sample $R^2$ nets off first-stage estimation error, so in the notation of \Cref{sec:theory} it behaves like $\lambda - s$ with $s = \sest/V$, and requiring $\lps \geq 0.3$ therefore requires $\lambda \geq 0.3 + s$. This makes the rule \emph{more} conservative toward instance normalization than the nominal threshold suggests, but it also means the anchor is a motivation for the order of magnitude of $\tau$ and not an identity between two coordinates. Second, the threshold is calibrated only in the sense of having been fixed before the data were seen; as \Cref{sec:empirical} reports, our eight datasets are bimodal in $\lps$ with nothing between $0.283$ and $0.744$, so the panel cannot distinguish a sharp crossover at $0.3$ from any cut placed inside that empty interval. Sensitivity of the rule to $\tau$ is quantified in \Cref{tab:tau} and should be read as a robustness check, not as evidence of fine calibration.

The three distinct crossover numbers that appear in this paper, and which of them anchors $\tau$, are set out in \Cref{app:thresholds}.


\subsection{Pre-registration as a falsifiability device}

Because both the diagnostic and the experimental grid are under our control, we bind our hands in advance. Before any grid run was executed, the official $\lps$ value of every dataset and the implied sign of the RevIN$-$CondNorm test-MSE gap were committed to the repository (git commit \texttt{cab17c1}), together with the success criterion (at least $6/8$ sign hits) and the fixed grid design. The pre-registration also flagged, in advance, electricity ($\lps = 0.283$, just below $\tau$) as the cell with the least confident sign prediction. All hits and misses are reported as-is in \Cref{sec:empirical}; the diagnostic is falsifiable, and \Cref{fig:lpsgap} shows the resulting predictions against the realized gaps.

\subsection{A known blind spot}

Absolute $\lps$ has a blind spot we did not resolve: on strongly persistent series (electricity market prices are the canonical example), covariates that are merely \emph{redundant} with the recent history can inflate $\lps$ even though they add nothing beyond what instance normalization already exploits. A natural incremental variant nets off a persistence baseline,
\begin{equation}\label{eq:dlps}
\dlps \;=\; R^2_{\mathrm{oos}}\bigl(\bar{y}_{i+1} \sim \bar{x}_{i+1},\, \bar{y}_i\bigr) \;-\; R^2_{\mathrm{oos}}\bigl(\bar{y}_{i+1} \sim \bar{y}_i\bigr),
\end{equation}
and we report it alongside $\lps$ and the persistence-only $R^2$ in \Cref{tab:tau} for completeness. We make no claim for it: on our eight datasets $\dlps$ never overturns an absolute-$\lps$ verdict, so the panel contains no case that would discriminate between the two, and validating an incremental rule requires the high-persistence, covariate-driven quadrant that our data do not sample. The pre-registered rule is absolute $\lps$ throughout.

\subsection{Why a pre-training diagnostic}

Choosing a normalization by grid search is expensive: our own comparison required 4{,}202 training runs across normalizations, backbones, horizons, and seeds. The $\lps$ computation, by contrast, involves only a handful of window-level first-stage fits---seconds to minutes on a CPU---and requires no forecasting model at all. Beyond cost, a pre-training diagnostic changes the epistemic status of the comparison: rather than selecting a normalization post hoc from test results, the practitioner commits to a choice from data properties alone. This is the move that feature-based forecasting-method selection made for whole methods \citep{talagala2023metalearning, monteromanso2020fforma}, applied here to a single pipeline component, and it answers directly the call of \citet{noise2025signal} for a principled diagnostics-driven strategy to replace the hand-set heuristic whose failure they report.
